% Primitive: derivative-toolkit — authored from calculus/derivatives_manim.py
% (six scenes) and its README concepts table; every number traces to plan
% 009's verified anchors (forward quotients throughout — the symmetric
% quotient of a quadratic is exact and kills the settling narrative).
% Backward-grounded: the zoom device, the mystery constants and the undo
% trick are owned by the e-and-ln chapter; the likelihood curve and the
% LSE ruler by the softmax chapter. The score function replaces the
% product rule everywhere (plan 009's keystone decision). Problem answers
% are computed by answers/derivative_toolkit.py (the solve-gate anchor).
\chapter{The derivative toolkit}

\begin{narrative}
The e-and-ln chapter built a device and never named its output: zoom
in on a smooth curve until it is straight, read the slope, and watch
the ratio settle as the step shrinks. This chapter names that ratio
$\frac{d}{dx}$ and packs the smallest toolkit the road needs — the sum
rule, the chain rule, two derivatives, and one reader called the
score. The kit is deliberately tiny for a reason the logarithms
chapter already gave: under the log, every product this road carries
becomes a sum, so almost nothing else is required.

\section{The slope is a function}

Run the zoom at several points of one parabola $f(x) = x^2$ and plot
the read-offs beneath it: every smooth curve carries a \emph{second}
curve, its slope at each point. At $x = 1$ the forward quotients
$\bigl((1+h)^2 - 1\bigr)/h$ for $h = 1, \tfrac{1}{10}, \tfrac{1}{100},
\tfrac{1}{1000}$ read
\[
  3, \quad 2.1, \quad 2.01, \quad 2.001 \;\longrightarrow\; 2,
\]
and no rounding is hiding anywhere: the quotient is exactly $2x + h$,
so the entries are literally $2 + h$, and ``settling'' means the $h$
fading out of the answer. The dual graph keeps two numbers visibly
different — at $x = 1$ the \emph{height} is $1$ while the
\emph{slope} is $2$ — the confusion a derivative most often drowns
in. One honest boundary: the device needs smooth. $|x|$ never
straightens at $0$, no matter the zoom. The notation is
Leibniz's~\cite{calculus-jeff-miller-earliest-uses-of-symbols-of-calculus} —
$dy/dx$ in a manuscript of November 1675, in print by
1684~\cite{calculus-wikipedia-nova-methodus-pro-maximis-et-minimis} —
and this book commits to it because it makes the chain rule look like
cancelling fractions, a promise examined two sections from now.

\section{Nudge in, nudge out}

For $x^2$ the derivative is drawn, not
computed~\cite{calculus-grant-sanderson-power-rule-through-geometry}.
Draw the square of side $x$ and grow the side by a nudge $dx$: the new
area is the old square, plus two strips of size $x \cdot dx$, plus a
corner of size $dx \cdot dx$. Halve the nudge and the strips halve —
but the corner \emph{quarters}: it is second-order small, and
discarding it is honest bookkeeping, not a
trick~\cite{calculus-silvanus-p-thompson-calculus-made-easy-1910}. At
$x = 3$ the strips say the slope is $6$; the finite check
$\bigl(3.01^2 - 9\bigr)/0.01 = 6.01$ agrees to the corner. The same
fact wears a second costume: near $x = 1$ the map $x \mapsto x^2$
stretches the number line by $\times 2$, near $x = 3$ by $\times 6$ —
a derivative is a \emph{local stretch factor}, the seed the chain
rule is about to harvest. Formula last:
\[
  d(x^2) \;=\; 2x\,dx + dx^2 ,
\]
with the $dx^2$ dying for a visible reason.

\section{Nudges add, nudges compose}

Two rules carry the whole kit. \emph{Changes add across a sum}:
heights stack, so their changes stack —
$d(f + g) = df + dg$, and constants ride along. And \emph{rates
multiply along a chain}: a nudge $dx$ causes a nudge $du$, which
causes a nudge $dy$, and the stretch factors compose,
\[
  \frac{dy}{dx} \;=\; \frac{dy}{du} \cdot \frac{du}{dx} .
\]
Worked honestly on $y = (2x)^2$ at $x = 1$: the inner map doubles
(rate $2$), the outer map squares — at the inner \emph{output}
$u = 2$, rate $2u = 4$ — so the chain predicts $2 \times 4 = 8$. The
forward quotients read $8.4$, $8.04$, $8.004 \to 8$. The table also
refutes the classic mistake on contact: evaluating the outer rate at
the input $x = 1$ instead of at $u = 2$ predicts $4$, and $4$ is not
where the quotients are headed. The $du$ cancellation is Leibniz's
notation keeping its promise — but the number lines are the
mechanism; the cancelling fraction is the mnemonic.

\section{The curve that is its own slope}

The e-and-ln chapter measured its mystery constants: $2^x$ grows at
$0.6931$ times its own height, $10^x$ at $2.3026$ times, and the
constants were the strides $\ln b$. In the new notation those
measurements were derivatives all along:
$\frac{d}{dx} b^x = \ln b \cdot b^x$, and $e$ is the base whose
constant is exactly $1$ — by definition and by measurement, nothing
proved twice. The ratio table $\bigl(e^{dt} - 1\bigr)/dt$ settles
through $1.7183$, $1.0517$, $1.0050$, $1.0005 \to 1$, so
\[
  \frac{d}{dx} e^x = e^x,
  \qquad
  \frac{d}{dx} \ln x = \frac{1}{x} .
\]
The second is one chain-rule line through the undo trick: differentiate
$e^{\ln x} = x$ to get $e^{\ln x} \cdot (\ln x)' = 1$, so
$(\ln x)' = 1/x$ — the debt-repaid flip, differentiated. The mirror
tells the same story in pictures: reflecting across $y = x$ swaps
rise and run, so slopes reciprocate — at $x = e$ the slope of
$\ln$ is $1/e$, mirroring $e^x$'s slope $e$ at height $e$. The
attribution is quotable: Euler named $e$ in 1748 (\emph{Introductio}
\S122)~\cite{calculus-leonhard-euler-introductio-in-analysin-infinitorum-i-vii-bru}
and wrote $d(e^x) = e^x\,dx$ in 1755 (\emph{Institutiones}
\S188)~\cite{calculus-leonhard-euler-institutiones-calculi-differentialis-i-vi-bru}.

\section{Zero slope finds the peak}

The softmax chapter found $\hat{p} = 3/4$ by sweeping a curve. The
sweep is over. The tool is the \emph{score function} — the counting
strip differentiated: under $\ln$, products become sums, so the
derivative of a product becomes a sum of \emph{relative rates},
\[
  \frac{d}{dx} \ln f \;=\; \frac{f'}{f} ,
\]
Euler's own sentence in
\S181~\cite{calculus-leonhard-euler-institutiones-calculi-differentialis-i-vi-bru}.
On the owned likelihood $L(p) = 4p^3(1-p)$:
$\ln L = \ln 4 + 3 \ln p + \ln(1-p)$ differentiates term by term —
the sum rule, $\ln'$, and one chain-rule line for $\ln(1-p)$ — into
the score
\[
  \frac{d \ln L}{dp} \;=\; \frac{3}{p} - \frac{1}{1-p} .
\]
The score runs $+4$ at $p = 1/2$, $+0.9524$ at $0.7$, $0$ at $3/4$,
$-1.25$ at $0.8$ — flipping sign exactly at the peak — and setting
it to zero is a \emph{linear} equation, $3(1-p) = p$, giving
$\hat{p} = 3/4$. In general $k/p - (n-k)/(1-p) = 0$ solves to
$\hat{p} = k/n$: the claim the random-variables chapter made about
proportions is now a theorem.

\begin{center}
\begin{tikzpicture}[x=6.4cm, y=6.4cm]
  % L(p) = 4 p^3 (1 - p) with the sign-change ribbon beneath — the
  % owned likelihood curve; peak at p = 3/4, L = 27/64 (plan 009 E).
  \draw[Muted] (0, 0) -- (1.04, 0);
  \draw[Muted] (0, 0) -- (0, 0.5);
  \draw[CoolInk, thick, smooth, domain=0:1, samples=60]
    plot (\x, {4*\x*\x*\x*(1 - \x)});
  \fill[AccentInk] (0.75, 0.421875) circle (0.55mm);
  \node[above, color=AccentInk] at (0.75, 0.45)
    {\small $\hat{p} = 3/4$};
  \node[left, color=CoolInk] at (-0.03, 0.42) {\small $L(p)$};
  % The ribbon: slope sign along the axis, zero exactly at 3/4.
  \fill[GoodInk!25] (0, -0.135) rectangle (0.75, -0.065);
  \fill[Warm!30] (0.75, -0.135) rectangle (1, -0.065);
  \node at (0.36, -0.1) {\small $+$};
  \node at (0.895, -0.1) {\small $-$};
  \draw[AccentInk, thick] (0.75, -0.15) -- (0.75, -0.05);
  \node[below, color=AccentInk] at (0.75, -0.17) {\small $0$};
  \node[below] at (0.02, -0.17) {\small $p = 0$};
  \node[below] at (0.98, -0.17) {\small $1$};
\end{tikzpicture}
\end{center}

Two honesty beats close the scene. The direct route
$dL/dp = 4p^2(3 - 4p)$ finds the same zero at $3/4$ — plus a second
root at $p = 0$ that is a valley floor, not a peak; and $x^3$'s slope
touches zero at the origin with no peak at all. Zero slope is
necessary, never sufficient — the sign \emph{change} is the
certificate, and checking it is the habit that survives every tool
built on top of this one. And the score itself needs $f > 0$: true
for likelihoods away from the endpoints, where the log is defined.

\section{The smooth max's shares}

The closer differentiates the softmax chapter's own objects. Nudge
$z_1 = 2$ by $0.01$ and the smooth max $\mathrm{LSE}(2, 1, 0) =
\anchor{008.L.LSE}$ moves by $\anchor{009.G.LSEnudge}$ — which is
$0.6652 \times 0.01$: $e^2$'s share of the total. One chain rule and
one sum rule prove the pattern
general~\cite{calculus-lawrence-murray-gradients-of-softmax-and-logsumexp}:
\[
  \frac{\partial}{\partial z_i} \mathrm{LSE}(z)
  \;=\; \frac{e^{z_i}}{\sum_j e^{z_j}}
  \;=\; \mathrm{softmax}(z)_i .
\]
The sensitivities of the smooth max \emph{are} the softmax shares —
the bars of the last chapter reborn as a gradient read-out, summing
to $1$ because the shares always did. And the loss follows in one
line: with the correct class $a$ scoring $2$,
$\mathrm{NLL} = \mathrm{LSE}(z) - z_a$ differentiates to
\[
  \nabla\,\mathrm{NLL} \;=\; \anchor{009.G.nllgradient}
  \;=\; p - \text{one-hot} :
\]
the softmax shares with exactly $1$ subtracted from the truth's
entry. (Sign convention, spoken once: this is the gradient of the
\emph{loss} — descending it is ascending the likelihood.) As the
correct score falls behind, the slope $p_c - 1$ walks
$\anchor{009.H.saturation}$, pressing toward $-1$ — the softmax
chapter's ``roughly linear gap'', now a theorem: a badly wrong
confident answer buys loss at essentially one unit per unit of score.

\paragraph{Where this goes.} A slope is advice: it says which way is
downhill, and the next chapter builds the rule that takes the advice
step by step. Beyond it, the road's destination hands this closing
picture to every frame of CTC — the same $p$ minus a target, but with
the one-hot softened into how often the truth actually used each
cell. The subtraction stays; only the target changes.
\end{narrative}

\section*{Practice problems}

\begin{problem}
Write the forward quotient of $f(x) = x^2$ at $x = 1$ for a step
$h$, simplify it exactly, and evaluate it at $h = 1, \tfrac{1}{10},
\tfrac{1}{100}$. What does the simplified form say ``settling''
means?
\begin{hint}
Expand $(1+h)^2$ before dividing.
\end{hint}
\begin{solution}
$\bigl((1+h)^2 - 1\bigr)/h = (2h + h^2)/h = 2 + h$ exactly — no
limit machinery needed to see it. The values are $3$, $2.1$,
$2.01$: each entry \emph{is} $2 + h$, so settling is nothing more
than the step fading out of the answer, leaving the slope $2$. (A
curiosity worth one aside: the symmetric quotient
$\bigl(f(1+h) - f(1-h)\bigr)/2h$ equals exactly $2$ for every $h$ —
a special property of quadratics, which is why the honest settling
story is told with forward quotients.)
\end{solution}
\end{problem}

\begin{problem}
Differentiate $y = (2x)^2$ at $x = 1$ by the chain rule, naming both
rates and where each is evaluated. A classmate answers $4$ —
reconstruct their mistake, and refute it with the forward quotients
at $h = \tfrac{1}{10}$ and $\tfrac{1}{100}$.
\begin{hint}
The outer rate is $2u$ — at which value of $u$?
\end{hint}
\begin{solution}
Inner map $u = 2x$: rate $2$. Outer map $y = u^2$: rate $2u$
evaluated at the inner \emph{output} $u = 2$, so $4$. The chain
multiplies: $2 \times 4 = 8$. The classmate evaluated the outer rate
at the input $x = 1$, getting $2 \times 2 = 4$. The quotients decide
it: $y = 4x^2$, so
$\bigl(4(1+h)^2 - 4\bigr)/h = 8 + 4h$, reading $8.4$ and $8.04$ —
headed for $8$, not $4$.
\end{solution}
\end{problem}

\begin{problem}
Derive $\frac{d}{dx} \ln x = 1/x$ from $\frac{d}{dx} e^x = e^x$ using
the undo identity $e^{\ln x} = x$, and check it numerically at
$x = 2$ with the forward quotient at $h = \tfrac{1}{100}$.
\begin{hint}
Differentiate both sides of the identity; the chain rule does the
rest.
\end{hint}
\begin{solution}
Differentiating $e^{\ln x} = x$: the left side, by the chain rule,
is $e^{\ln x} \cdot (\ln x)'$, and the right side is $1$. But
$e^{\ln x} = x$, so $x \cdot (\ln x)' = 1$ and $(\ln x)' = 1/x$ —
the undo-never-cancel flip carrying its own derivative. At $x = 2$:
$\bigl(\ln 2.01 - \ln 2\bigr)/0.01 = 0.4988$, already close to $1/2$
and settling as $h$ shrinks (the quotients run $0.4879$, $0.4988$,
$0.4999 \to 0.5$ at $h = 10^{-1}, 10^{-2}, 10^{-3}$).
\end{solution}
\end{problem}

\begin{problem}
For the likelihood $L(p) = 4p^3(1-p)$: write $\ln L$, differentiate
it term by term into the score, evaluate the score at $p = 1/2$,
$0.7$, and $0.8$, and solve the zero. Then find the general
$\hat{p}$ for $k$ heads in $n$ flips.
\begin{hint}
$\ln(1-p)$ needs one chain-rule line; the zero is a linear equation.
\end{hint}
\begin{solution}
$\ln L = \ln 4 + 3 \ln p + \ln(1-p)$, so the score is
$3/p - 1/(1-p)$ (the constant contributes nothing; $\ln(1-p)$
differentiates to $-1/(1-p)$ by the chain rule). Values:
$6 - 2 = +4$ at $p = 1/2$; $3/0.7 - 1/0.3 = 30/7 - 10/3 = 20/21
\approx +0.9524$ at $0.7$; $15/4 - 5 = -5/4 = -1.25$ at $0.8$. The
sign flips across $p = 3/4$, and the zero is linear:
$3/p = 1/(1-p) \Rightarrow 3(1-p) = p \Rightarrow \hat{p} = 3/4$. In
general the score is $k/p - (n-k)/(1-p)$, whose zero
$k(1-p) = (n-k)p$ gives $\hat{p} = k/n$ — the observed proportion,
for every $k$ and $n$ at once.
\end{solution}
\end{problem}

\begin{problem}[stretch]
With the correct class scoring $2$ in $z = (2, 1, 0)$: differentiate
$\mathrm{NLL}(z) = \mathrm{LSE}(z) - z_a$ component by component,
and evaluate the gradient to four decimal places. Then let the
correct class's score fall: on $z = (2, 1, t)$ with the truth in the
third slot, compute the slope $\partial\,\mathrm{NLL}/\partial t$ at
$t = 0, -2, -5$. What is the slope pressing toward, and why can it
never arrive?
\begin{hint}
$\partial_i \mathrm{LSE} = \mathrm{softmax}(z)_i$, and $z_a$
contributes $-1$ to exactly one component. For the walk, the slope
is $p_c - 1$.
\end{hint}
\begin{solution}
Each component of $\nabla \mathrm{LSE}$ is a softmax share, and
subtracting $z_a$ subtracts $1$ from the truth's component:
$\nabla\,\mathrm{NLL} = p - \text{one-hot} =
(0.6652 - 1,\ 0.2447,\ 0.0900) = (-0.3348,\ 0.2447,\ 0.0900)$,
whose exact components sum to $0$ (the shares sum to $1$, and one
$1$ was removed). On $z = (2, 1, t)$ the slope in $t$ is
$p_c - 1$ where $p_c = \mathrm{softmax}(z)_3$: at $t = 0, -2, -5$ it
reads $-0.9100$, $-0.9868$, $-0.9993$. The walk presses toward $-1$
but never arrives: $p_c$ is a ratio of positive exponentials, so it
is positive at every finite $t$ — the loss becomes as straight as
you like, with slope $-1$ only as a limit. (At $t = -10$ the share
is $\approx 4.5 \times 10^{-6}$; write it in scientific notation —
rounding the slope to $-1.0000$ would claim the limit was reached.)
\end{solution}
\end{problem}
